\newsection{Object-Oriented Analysis, Design \& Architecture}{sec:ooad}

\subsection{Object-Oriented Analysis}
Object-oriented analysis finds the classes of the problem domain before any design decision is made. We used the noun analysis method [2]: mark the nouns of the problem statement, classify them, filter out the ones that are not classes, and add classes that domain knowledge requires.

\head{Step 1: Potential Objects (Nouns and Noun Phrases)}
From the problem statement in Section~\ref{sec:objectives}: department, system, calendar, faculty member, HOD, timetable, timetable file, cell, initials, teacher, class, term, holiday, photo, meeting circular, date, time, venue, name, meeting, details, reminder, e-mail, phone notification, quiet hours, documentation task, group, deadline, work, alert.

\head{Step 2: Classification}
\needspace{7\baselineskip}
\begin{xltabular}{\textwidth}{|p{3.2cm}|Y|}
\caption{Classification of candidate objects}\\
\hline
\textbf{Category} & \textbf{Candidates} \\ \hline
People and roles & faculty member, teacher, HOD, group \\ \hline
Organisations & department \\ \hline
Things and documents & timetable file, cell, photo, meeting circular, e-mail, phone notification \\ \hline
Concepts & calendar, initials, quiet hours, details, term \\ \hline
Events & class, meeting, holiday, deadline, reminder, alert, documentation task \\ \hline
Attributes & date, time, venue, name \\ \hline
\end{xltabular}

\head{Step 3: Filtering}
\needspace{7\baselineskip}
\begin{xltabular}{\textwidth}{|p{3.6cm}|p{3.4cm}|Y|}
\caption{Filtering of candidates}\\
\hline
\textbf{Candidate} & \textbf{Decision} & \textbf{Reason} \\ \hline
system & Removed & It is the system itself, not a domain object. \\ \hline
teacher, faculty member & Merged $\rightarrow$ FacultyProfile & Same person; profile holds initials and aliases. \\ \hline
HOD, Admin & Role of User & They differ only in permissions (\texttt{Role} value). \\ \hline
class, meeting, deadline & Merged $\rightarrow$ Event & One interval type with a \texttt{kind}; clash checks and reminders treat them alike. \\ \hline
reminder, alert & Merged $\rightarrow$ ScheduledNotification & One planned send with its own life cycle. \\ \hline
date, time, venue, name & Attributes & Properties of Event or FacultyProfile. \\ \hline
initials, quiet hours & Value objects & Defined only by their value; must validate themselves. \\ \hline
cell & Merged $\rightarrow$ TimetableSlot & A parsed cell is one weekly slot. \\ \hline
e-mail, phone notification & Channel value & Delivery channel of a notification. \\ \hline
calendar, details, group & Removed & A view over events, too vague, or a set of profiles. \\ \hline
\end{xltabular}

\head{Step 4: Classes Added from Domain Knowledge}
\begin{itemize}
  \item \textbf{User} -- the log-in account, separate from the faculty profile (an Admin may have no profile).
  \item \textbf{TimetableUpload} -- one uploaded file with its parse report and status.
  \item \textbf{EventAttendee} -- links a faculty member to an event with a response.
  \item \textbf{ReminderRule} -- a user's lead times and channels for each kind of event.
  \item \textbf{AcademicTerm} -- the term dates and time zone that bound class generation.
\end{itemize}

\head{Step 5: Final Domain Classes}
Department, User, FacultyProfile, AcademicTerm, Holiday, Course, Section, Room, TimetableUpload, TimetableSlot, Event, EventAttendee, Circular, ScheduledNotification, ReminderRule, PushSubscription, Task, TaskTemplate, AuditLog, RefreshToken; value objects Initials, TimeSlot and QuietHours.

\subsection{Object-Oriented Design}

\head{Layered (Clean) Architecture}
The backend is divided into four layers [5]. Dependencies point inwards, towards the domain.

\begin{figure}[!htbp]
\centering
\begin{tikzpicture}[
  lay/.style={rectangle, draw, semithick, align=center, font=\footnotesize, minimum width=12.5cm, minimum height=1.25cm, text width=12cm},
  ar/.style={semithick, dashed, -{Stealth[length=2.2mm]}}]
  \node[lay] (api) at (0,0) {\textbf{API layer} (\texttt{api/}) and \textbf{workers} (\texttt{workers/})\\FastAPI routers, request/response schemas, dependency wiring; Celery tasks};
  \node[lay] (app) at (0,-2.2) {\textbf{Application layer} (\texttt{application/})\\one class per use case; clash detection, recurrence, reminder and task algorithms};
  \node[lay] (dom) at (0,-4.4) {\textbf{Domain layer} (\texttt{domain/})\\entities, value objects, domain events, errors, repository interfaces --- no framework imports};
  \node[lay] (inf) at (0,-6.6) {\textbf{Infrastructure layer} (\texttt{infrastructure/})\\SQLAlchemy repositories, parsers, OCR, notifiers, MinIO storage, rate limiter, backup};
  \draw[ar] (api) -- node[right, font=\scriptsize] {uses} (app);
  \draw[ar] (app) -- node[right, font=\scriptsize] {uses} (dom);
  \draw[ar] (inf) -- node[right, font=\scriptsize] {implements interfaces of domain} (dom);
\end{tikzpicture}
\caption{Layered architecture of the backend}
\label{fig:layers}
\end{figure}

Use cases receive their collaborators through their constructors (dependency injection). The same \texttt{ProcessCircularUseCase} therefore runs in a worker in production and in a test with a test database, and the same \texttt{ReminderDispatcher} runs with \texttt{SystemClock} in production and with \texttt{FrozenClock} in tests. As shown in the package diagram, four application files still import infrastructure classes directly; the domain layer has no outward imports at all.

\head{Design Principles Applied}
\needspace{7\baselineskip}
\begin{xltabular}{\textwidth}{|p{3.6cm}|Y|}
\caption{Design principles}\\
\hline
\textbf{Principle} & \textbf{How FacultyDesk applies it} \\ \hline
Single responsibility & Each use-case class serves one goal; \texttt{InitialsResolver} only resolves initials; \texttt{CircularExtractor} only extracts fields. \\ \hline
Open--closed & A new timetable format needs a new \texttt{TimetableParser} subclass and one line in the factory; a new channel (such as SMS) needs a new \texttt{Notifier}. \\ \hline
Liskov substitution & Any \texttt{Notifier} can be given to the dispatcher; tests substitute frozen-time and test-database objects. \\ \hline
Interface segregation & Eighteen small repository interfaces, one per aggregate, instead of one large data-access class. \\ \hline
Dependency inversion & Use cases depend on \texttt{UnitOfWork}, \texttt{Notifier} and \texttt{Clock} abstractions, not on SQLAlchemy or SMTP. \\ \hline
Information expert (GRASP) & \texttt{QuietHours} decides whether a time is quiet; \texttt{Circular} decides which status change is legal; \texttt{Task} validates its own importance and deadline. \\ \hline
\end{xltabular}

\head{Design Patterns}
The patterns follow the catalogues of Gamma \emph{et al.} [3] and Fowler [4]. Figures~\ref{fig:classpatterns} and \ref{fig:classpatternsb} show the classes involved.

{\small
\needspace{7\baselineskip}
\begin{xltabular}{\textwidth}{|p{2.6cm}|p{4.6cm}|Y|}
\caption{Design patterns used}\\
\hline
\textbf{Pattern} & \textbf{Classes} & \textbf{Problem it solves} \\ \hline
Repository & EventRepository, TaskRepository and 16 others & Hides SQL behind collection-like interfaces so domain code stays storage-free. \\ \hline
Unit of Work & UnitOfWork, SqlAlchemyUnitOfWork & Saves a circular, its meeting, attendees and reminders in one transaction. \\ \hline
Strategy & TimetableParser subclasses; Notifier subclasses & Swaps the parsing algorithm by file type and the delivery method by channel. \\ \hline
Factory & TimetableParserFactory & Chooses the parser from the file's first bytes, not from its name. \\ \hline
Composite & CompositeNotifier & Lets the dispatcher send through e-mail and push as if they were one notifier. \\ \hline
State & Circular, ScheduledNotification, TimetableUpload & Encodes legal life cycles as transition tables; illegal moves raise errors. \\ \hline
Domain events & EventScheduled, EventRescheduled, EventCancelled, TaskDeadlineChanged, TimetableConfirmed & Records what happened to an entity; reminder re-planning reacts to rescheduling and cancellation. \\ \hline
\end{xltabular}}

\begin{figure}[!htbp]
  \centering
  \resizebox{\textwidth}{!}{\begin{tikzpicture}[
    inh/.style={semithick, -{Triangle[open, length=3mm, width=3mm]}},
    dep/.style={semithick, dashed, -{Stealth[length=2.2mm]}},
    agg/.style={semithick, {Diamond[open, length=3mm, width=2mm]}-{Stealth[length=2.2mm]}},
    lbl/.style={font=\scriptsize\itshape, fill=white, inner sep=1pt}
  ]
  \umlclass{factory}{0,3.4}{TimetableParserFactory}{}{
    + get\_parser(content, filename)\\
    \phantom{+} : TimetableParser}
  \umlclass{parser}{0,0}{<{}<abstract>{}>\\TimetableParser}{}{
    + parse(content) : ParseResult}
  \umlclass{xlsx}{-4.4,-3.4}{XlsxTimetableParser}{}{+ parse(content)}
  \umlclass{csv}{0,-3.4}{CsvTimetableParser}{}{+ parse(content)}
  \umlclass{pdf}{4.4,-3.4}{PdfTimetableParser}{}{+ parse(content)}
  \draw[dep] (factory) -- node[lbl, right] {<{}<creates>{}>} (parser);
  \draw[inh] (xlsx.north) -- ++(0,0.6) -| (parser.south);
  \draw[inh] (csv.north) -- (parser.south);
  \draw[inh] (pdf.north) -- ++(0,0.6) -| (parser.south);
  \node[font=\footnotesize\bfseries, anchor=south] at (0,5.0) {Strategy + Factory};

  \umlclass{composite}{11.4,3.4}{CompositeNotifier}{
    - email : EmailNotifier\\
    - push : WebPushNotifier}{+ send(...)}
  \umlclass{notifier}{11.4,0}{<{}<abstract>{}>\\Notifier}{}{
    + send(notification, payload, ...)}
  \umlclass{email}{9.2,-3.4}{EmailNotifier}{
    - hostname, port : SMTP}{+ send(...)}
  \umlclass{push}{13.6,-3.4}{WebPushNotifier}{
    - vapid\_private\_key : str}{+ send(...)}
  \draw[inh] (composite.south) -- (notifier.north);
  \draw[agg] (composite.east) -- ++(0.7,0) |- node[lbl, pos=0.25, right] {2} (notifier.east);
  \draw[inh] (email.north) -- ++(0,0.6) -| (notifier.south);
  \draw[inh] (push.north) -- ++(0,0.6) -| (notifier.south);
  \node[font=\footnotesize\bfseries, anchor=south] at (11.4,5.0) {Strategy + Composite};
\end{tikzpicture}
}
  \caption{Class diagram of the Strategy, Factory and Composite patterns}
  \label{fig:classpatterns}
\end{figure}

\begin{figure}[!htbp]
  \centering
  \resizebox{\textwidth}{!}{\begin{tikzpicture}[
    inh/.style={semithick, -{Triangle[open, length=3mm, width=3mm]}},
    real/.style={semithick, dashed, -{Triangle[open, length=3mm, width=3mm]}},
    dep/.style={semithick, dashed, -{Stealth[length=2.2mm]}},
    lbl/.style={font=\scriptsize\itshape, fill=white, inner sep=1pt}
  ]
  \umlclass{uow}{0,0}{<{}<abstract>{}>\\UnitOfWork}{
    + events : EventRepository\\
    + tasks : TaskRepository\\
    + circulars : CircularRepository\\
    + notifications : NotificationRepository}{
    + commit() / rollback()}
  \umlclass{sqlauow}{0,-4.6}{SqlAlchemyUnitOfWork}{
    - session : AsyncSession}{+ commit() / rollback()}
  \draw[inh] (sqlauow) -- (uow);
  \umlclass{disp}{5.8,0}{ReminderDispatcher}{
    - clock : Clock\\
    - notifier : Notifier}{
    + dispatch\_due(uow, batch) : int\\
    + recover\_stale\_sending(uow) : int}
  \umlclass{notif}{5.8,-4.6}{<{}<abstract>{}>\\Notifier}{}{+ send(...)}
  \umlclass{clock}{11.6,0}{<{}<protocol>{}>\\Clock}{}{+ now() : datetime}
  \umlclass{sysclock}{9.6,-4.6}{SystemClock}{}{+ now()}
  \umlclass{frozen}{13.6,-4.6}{FrozenClock}{}{+ now()}
  \draw[real] (sysclock.north) -- ++(0,0.9) -| (clock.south);
  \draw[real] (frozen.north) -- ++(0,0.9) -| (clock.south);
  \draw[dep] (disp) -- node[lbl, above] {uses} (uow);
  \draw[dep] (disp) -- node[lbl, above] {uses} (clock);
  \draw[dep] (disp) -- node[lbl, right] {uses} (notif);
  \node[font=\footnotesize\bfseries, anchor=south] at (0,2.1) {Repository + Unit of Work};
  \node[font=\footnotesize\bfseries, anchor=south] at (11.6,2.1) {Dependency injection of Clock};
\end{tikzpicture}
}
  \caption{Class diagram of the Repository and Unit of Work patterns and the injected Clock}
  \label{fig:classpatternsb}
\end{figure}

\subsection{Architecture Decisions}
\needspace{7\baselineskip}
\begin{xltabular}{\textwidth}{|p{3.4cm}|Y|}
\caption{Architecture decisions}\\
\hline
\textbf{Decision} & \textbf{Reason} \\ \hline
Background workers for slow work & OCR and timetable parsing take seconds; the API stores the file, replies \texttt{202 Accepted} and a Celery worker does the work, so every web request stays short. \\ \hline
UTC storage, local expansion & Every timestamp is stored in UTC. Weekly recurrence is expanded in local time (Asia/Kolkata) and then converted, so a 9:00 class stays at 9:00. \\ \hline
Injected clock & No code reads the system time directly; an injected \texttt{Clock} makes every time rule testable. \\ \hline
Exactly-once reminders & Key $k = \mathrm{SHA256}(\text{target}\,\|\,\text{user}\,\|\,\text{offset}\,\|\,\text{channel})$ with a unique database constraint; rows claimed with \texttt{FOR UPDATE SKIP LOCKED}. \\ \hline
Human in the loop & No OCR result reaches a calendar unless it passes the review gate and auto-confirmation is enabled; otherwise the HOD confirms it. \\ \hline
One authorisation policy & One \texttt{AuthorisationPolicy} class decides what Admin, HOD and Faculty may do, so permission rules live in one place. \\ \hline
\end{xltabular}
